% (C) 2026 Ralf Stephan, in collaboration with Claude Code.
% Released under CC0 1.0 Universal (public-domain dedication).
% See https://creativecommons.org/publicdomain/zero/1.0/

\documentclass[11pt]{article}
\usepackage[T1]{fontenc}
\usepackage[utf8]{inputenc}
\usepackage{lmodern}
\usepackage[margin=2.5cm]{geometry}
\usepackage{amsmath,amssymb}
\usepackage{braket}
\usepackage[hidelinks]{hyperref}

% Lean names and paths: verbatim, breakable at . _ /
\DeclareUrlCommand\lean{\urlstyle{tt}}
\def\UrlBreaks{\do\.\do\_\do\/}
% Marks an open problem that blocks checking a main claim of the brief
\newcommand{\blocker}[1]{\textsc{[blocks verification: #1]}~}
% Extra stretch for paragraphs with long, unbreakable Lean names
\emergencystretch=6em

\title{Blockers:\\ Unitary evolution as infinitesimal projective processes}
\author{}
\date{}

\begin{document}
\maketitle

\noindent
This file collects everything about the two brief issues of \lean{open-problems.tex} that block
verification of Theorem~18, together with a proposed resolution of each and the changes it
needs in the brief. Sections~\ref{sec:blockers} and~\ref{sec:related} are copied from
\lean{open-problems.tex} and \lean{formalization-log.tex}; the rest is new.

%======================================================================
\section{The blockers}\label{sec:blockers}

From \lean{open-problems.tex}, ``From the formalization'':
\begin{enumerate}
\item \blocker{brief issue}
  \textbf{Durations with irrational ratio.}
  DR-PROJ defines $\Phi_t$ only as
  $\lim_n(\Psi_{t/n})^n$. For $t/s$ irrational a common step $\tau$ gives blocks $\Psi_\tau^k$ that
  are not of this form, so ``common refinements'' and continuity at zero do not prove
  $\Phi_{t+s} = \Phi_t\circ\Phi_s$; eq.~(26) would have to hold uniformly along the trajectory.
  The remaining interval itself is harmless. A fix: in finite dimension a telescoping estimate
  gives $\Phi_t = e^{tL}$ directly, which is formalized (\lean{IsProjectiveLimit.exists_eq_exp});
  in infinite dimension Chernoff's product formula does, if the closure of $L$ generates a
  contraction semigroup, which is open. In infinite dimension this is the one clause of
  \lean{IsEvolution} that is assumed; the others, including composition for durations of
  rational ratio (eq.~(29)), are derived (\lean{IsProjectiveLimit.isEvolution}). $\Phi_0 = I$ is
  assumed in both dimensions, since the formalization defines $\Phi_t$ only for $t > 0$; the brief
  obtains it from $\Psi_0 = I$.

  The gap may be real in infinite dimension. A sketch, not checked in Lean: let $H_1 \ne H_2$ be
  two self-adjoint extensions of a symmetric operator $H_0$ with deficiency indices $(1,1)$, and let
  $\Psi_\tau$ be the cutoff resolvent of $H_1$ for rational $\tau$ and of $H_2$ for irrational
  $\tau$. Since $t/n$ is rational exactly when $t$ is, limits, prediction and reconstruction hold
  for each $t$, and $\Phi_t$ is conjugation by $e^{-itH_1}$ or $e^{-itH_2}$. If both families have
  the derivative $[H_0,\rho]/(i\hbar)$ on the density operators built from vectors of
  $\mathcal{D}(H_0)$, which are dense, then DR-PROJ holds while DR-UNIT would force
  $H_1 = H_2 + c$. (Brief, Sec.~4, after eq.~(29); Log, ``Problems found''.)
\item \blocker{brief issue}\textbf{What perfect reconstruction means.}
  DR-PROJ only asks that the probabilities of
  the reconstructed sequences tend to one. The formalization instead assumes a second projective
  limit $\Phi'_t$ with $\Phi'_t(\Phi_t P) = P$ and $\Phi_t(\Phi'_t P) = P$ on pure states. It is
  not settled whether the brief's weaker wording already makes $\Phi_t$ a bijection of the pure
  states, or what the right formulation is. (Brief, Sec.~4: ``Perfect reconstruction gives the
  converse correspondence''; Log, ``Differences'' and ``Problems found''.)
\end{enumerate}

%======================================================================
\section{Related items}\label{sec:related}

From \lean{open-problems.tex}, ``From the formalization'' (the preamble, which places the
blockers in the any-dimension formalization):

\begin{quote}
Beyond finite dimension the formalization has a part in any dimension, which covers
unbounded $H$ in DR-UNIT $\Rightarrow$ DR-PROJ and derives the evolution structure
\lean{IsEvolution} from DR-PROJ up to the item ``Durations with irrational ratio'' below
(\lean{drUnit_iff_drProj_and_map_zero_and_map_add}), for every cutoff level satisfying
eq.~(22). It leaves open the item ``Resolvent of an unbounded generator'' above.
\end{quote}

\noindent
The former item ``The generator of the cutoff family'' of \lean{open-problems.tex}, needed if the
first blocker is resolved through the generator by Chernoff's formula, is resolved. On the core,
the generator
of the family~(23) is $\mathcal{L}(\rho) = [H,\rho]/(i\hbar)$: $\rho$ maps into the domain of $H$,
and $\mathcal{L}(\rho)\psi = -i(H\rho\psi - \rho H\psi)$ for $\psi$ in the domain
(\lean{cutoffResolvent_generator}; Log, Sec.~1.1, ``Differences'', item on eq.~(26)).

\noindent
From \lean{open-problems.tex}, ``From the brief'' (likewise):
\begin{enumerate}
\setcounter{enumi}{1}
\item \textbf{Resolvent of an unbounded generator.} The brief states that ``direct use of the
  resolvent of an unbounded generator still gives a completely positive trace-preserving map
  and the resolvent exponential formula'', and that the spectral truncations are needed only
  for states outside the domain of $H$. Neither claim is proved, and whether the untruncated
  $(I - \tau\mathcal{L})^{-1}$ gives perfect prediction on the domain of $H$ is left open.
  Not formalized. (Brief, Sec.~5, ``Unbounded generators''; Log, Sec.~1.1, ``Differences'',
  last item.)
\end{enumerate}

\noindent
From \lean{formalization-log.tex}, Section~1, ``Differences from the brief'':
\begin{itemize}
\item DR-PROJ is made precise as follows. Outcomes and branches refer to a projective
  decomposition that realizes $\Psi$ on the current pure state. Perfect prediction asks that the
  final branch state converges, without naming the limit; DR-PROJ $\Rightarrow$ DR-UNIT shows that
  it is $\Phi_tP$. Reconstruction is a second projective limit that inverts $\Phi_t$ on pure
  states. For the resolvents it is the limit of $(1 + \tau L)^{-1}$, as in the remark ``The
  finite process is not unitary''.
\item \lean{drUnit_iff_drProj} assumes \lean{IsEvolution} of the brief on the Schr\"odinger
  equation: $\Phi_t$ is affine and preserves density operators, $\Phi_0 = I$,
  $\Phi_{s+t} = \Phi_s\Phi_t$, and $t \mapsto \Phi_t\rho$ is continuous at $0^+$. The brief
  derives these from DR-PROJ (affinity of the limit, eqs.~(28)--(29), common refinements).
  \lean{IsProjectiveLimit.isEvolution} derives them too, but not by the brief's route. Each
  $\Psi_\tau$ agrees on density operators with a linear map $T_\tau$, because the density
  operators contain a basis of the operators; $(T_\tau - I)/\tau \to B$ in operator norm, because
  convergence on a basis is norm convergence; and $(T_{t/n})^n \to e^{tB}$ by telescoping. So
  $\Phi_t = e^{tB}$ on density operators for every $t > 0$, which gives affinity, preservation of
  density operators, continuity at $0^+$ and composition for all durations at once; eq.~(28) and
  the common refinements are not needed. DR-PROJ says nothing about $\Phi_0$, which is assumed to
  be the identity: \lean{drUnit_iff_drProj_and_map_zero} is DR-UNIT $\Leftrightarrow$ DR-PROJ
  with $\Phi_0 = I$.
\end{itemize}

\noindent
From \lean{formalization-log.tex}, Section~1, ``Problems found in the brief'':
\begin{itemize}
\item Section~3 states DR-PROJ only as ``the probabilities of the predicted and reconstructed
  state sequences tend to one''. The remark on prediction and reconstruction adds that the
  predicted states converge. This is not an extra assumption: since the branch is within
  $2(1-w)$ of $\Psi^nP \to \Phi_tP$, the predicted states converge to $\Phi_tP$ anyway. The
  finite-dimensional proof still uses the clause; the proof in any dimension does not
  (Section~1.1).
\item ``Perfect reconstruction gives the converse correspondence'' does not say how the
  reconstructed sequence relates to the predicted one. Injectivity and surjectivity of $\Phi_t$ on
  pure states need the reconstruction to invert $\Phi_t$; this is what DR-PROJ assumes here.
\item Composition for durations with irrational ratio, ``take common refinements with a
  remaining interval whose duration tends to zero'', is a gap. The remaining interval is
  harmless: with a common step $\tau$ the powers of $\Psi_\tau$ commute, so
  $\|\Psi_\tau^k(\Psi_\tau^j\rho - \rho)\|_1 \le j\|\Psi_\tau\rho - \rho\|_1
  \approx j\tau\|L\rho\|_1
  \to 0$ by the telescoping of eq.~(28). The problem is the step size. DR-PROJ defines $\Phi_t$
  only as $\lim_n (\Psi_{t/n})^n$. With $t/s$ irrational and $\tau = (t+s)/N$, the block of
  $k = \lfloor t/\tau\rfloor$ steps is $\Psi_\tau^k$, which is not of the form $(\Psi_{t/k})^k$,
  and DR-PROJ says nothing about its limit. Comparing $(\Psi_a)^k$ with $(\Psi_b)^k$ for nearby
  $a, b$ needs the $o(\tau)$ in eq.~(26) to be uniform along the states $\Psi_b^j\rho$, while
  eq.~(26) holds for each $\rho$ separately. Continuity at zero does not supply this. In finite
  dimension $L$ is bounded, eq.~(26) holds in operator norm, and telescoping gives
  $\|(\Psi_{t/n})^n - e^{tL}\| \le n\,o(t/n)\,e^{t\|L\|} \to 0$, so $\Phi_t = e^{tL}$ and
  composition follow at once (formalized, \lean{IsProjectiveLimit.exists_eq_exp}). In infinite
  dimension Chernoff's product formula gives the same,
  uniformly on compact $t$-intervals, provided the closure of $L$ generates a contraction
  semigroup; the brief's ``single, possibly unbounded generator $L$'' must be read to include
  this, or it must be assumed.
\end{itemize}

\noindent
From \lean{formalization-log.tex}, Section~1.1 (any dimension), ``Differences from the brief'':
\begin{itemize}
\item DR-PROJ $\Rightarrow$ DR-UNIT is proved over \lean{IsEvolution} (\lean{DrProj.drUnit}) and
  over two of its clauses only (\lean{drUnit_iff_drProj_and_map_zero_and_map_add}). The
  finite-dimensional derivation of \lean{IsEvolution}, through $\Phi_t = e^{tB}$, needs a finite
  basis of density operators; \lean{IsProjectiveLimit.isEvolution} follows the brief instead.
  Each $\Phi_t$ is a trace-norm limit of affine maps that keep density operators, so it is one
  too, and contracts the trace norm (Simon's lemma, the brief's ``extend to positive
  trace-preserving linear maps''). Eq.~(28) is the telescoping bound
  $\|\Psi^n\rho - \rho\|_1 \le n\|\Psi\rho - \rho\|_1$ on the dense set of eq.~(26), and
  contractivity extends continuity at $0^+$ to every density operator. Eq.~(29) is proved as
  $\Phi_{(p+q)u} = \Phi_{pu}\circ\Phi_{qu}$, along $n = (p+q)m$, which gives composition whenever
  $s/t$ is rational. Two clauses remain assumptions: $\Phi_0 = I$, since \lean{IsProjectiveLimit}
  defines $\Phi_t$ only for $t > 0$ (the brief obtains it from $\Psi_0 = I$), and
  $\Phi_{s+t} = \Phi_s\circ\Phi_t$ for $s/t$ irrational, the gap noted in the finite-dimensional
  part.
\item The proof of DR-PROJ $\Rightarrow$ DR-UNIT does not use the clause of
  \lean{PerfectPrediction} that the final branch states converge. The bound
  $\|\Psi^nP - Q_n\|_1 \le 2(1-w_n)$ and $\Psi^nP \to \Phi_tP$ already give $Q_n \to \Phi_tP$.
  The clause is kept so that the definitions match those of finite dimension.
\end{itemize}

\noindent
From \lean{formalization-log.tex}, Section~1.1, ``Problems found in the brief'':
\begin{itemize}
\item Composition for durations of irrational ratio is the gap noted in the finite-dimensional
  part. Here no telescoping estimate closes it: $L$ may be unbounded, and eq.~(26) holds only
  for each $\rho$ of a dense set, not uniformly.
\end{itemize}

%======================================================================
\section{The author's proposal (2026-11-02)}

Paraphrased, from a reply on another brief that the author means for all briefs: the briefs get a
section of mathematical prerequisites, which states what Lean should prove and is fully
formalized. The relations between physical conditions, such as Theorem~18, are argued informally
from these prerequisites and are checked by reading, not by Lean. The author will rewrite the
briefs accordingly.

This brief already has such a section (Sec.~2: Definition~1, Propositions~2, 13 and~17), and it
is formalized. Both blockers sit in Theorem~18, in what DR-PROJ means. Neither is a question of
reading a physical condition charitably: each decides which mathematical statement the proof of
Theorem~18 rests on. So DR-PROJ has to be stated so that the proof goes through, and the facts
the proof then uses belong in the prerequisites.

%======================================================================
\section{Proposed resolution}

\subsection{Durations with irrational ratio}

\paragraph{Guess.} DR-PROJ means that the evolution does not depend on how the elapsed time is
subdivided. The limit is taken over \emph{every} subdivision of $[0,t]$ whose longest interval
tends to zero, not only over the equal subdivisions $t/n$:
\[
  \Phi_t\rho = \lim \Psi_{\tau_{k}} \circ \dots \circ \Psi_{\tau_{1}}\,\rho
  \qquad\text{in trace norm, for } \tau_i > 0,\ \textstyle\sum_i \tau_i = t,\
  \max_i \tau_i \to 0,
\]
and $\Phi_0 = I$, the empty product. Composition is then immediate (prerequisite P1): a
subdivision of $[0,s]$ followed by one of $[0,t]$ is a subdivision of $[0,s+t]$. DR-UNIT still
implies the stronger DR-PROJ (P2), by the telescoping argument of the brief, run on the core of
bounded spectral support. The equal subdivisions are a special case, so every other step of the
proof stands as it is. The counterexample sketch of Section~\ref{sec:blockers} is excluded: a
subdivision of a rational $t$ into irrational steps leads to $e^{-itH_2}$, not $e^{-itH_1}$.

\paragraph{Alternative.} Keep the equal subdivisions and assume instead that the closure of the
generator $L$ of eq.~(26) generates a contraction semigroup. Chernoff's product formula then gives
$\Phi_t = e^{t\bar L}$, and composition follows. For the direction DR-UNIT $\Rightarrow$ DR-PROJ
this needs $\bar L$ to be the generator of $\rho \mapsto U_t\rho U_t^\dagger$; on the core, $L$ is
$[H,\cdot\,]/(i\hbar)$ (\lean{cutoffResolvent_generator}), but that the core is a core for this
generator is not formalized, and Chernoff's formula is not in Mathlib.

\paragraph{Why.}
\begin{itemize}
\item It is what the brief already says. ``The use of the same duration-dependent process makes
  the limiting evolution depend only on elapsed time'' and ``Composition follows by taking common
  refinements'' make sense only if unequal steps are allowed: a common refinement of $t$ and $s$
  is not an equal subdivision of $t+s$ when $t/s$ is irrational.
\item Time is homogeneous. An observer who checks the system at uneven times, say on a clock that
  ticks irregularly, should see the same evolution as one who checks at even times. A process for
  which the outcome depends on whether the step lengths are rational is not a physical process.
\item It costs no physics. The direction DR-UNIT $\Rightarrow$ DR-PROJ survives with the same
  estimates (P2), so Theorem~18 keeps its full range, including unbounded generators.
\item It is elementary. P1 is a three-line argument from trace-norm contraction, and P2 reuses
  the second-order bounds of Proposition~13. The alternative needs semigroup theory that neither
  the brief nor Mathlib has.
\end{itemize}

\paragraph{Confidence.}
\begin{itemize}
\item High: some strengthening of DR-PROJ is needed. The counterexample sketch is not checked in
  Lean, but no argument from equal subdivisions alone is known.
\item Medium-high: the subdivision form is what the author means, given ``common refinements''.
  The author may prefer the generator form; both resolve the blocker.
\end{itemize}

\subsection{What perfect reconstruction means}

\paragraph{Guess.} Reconstruction is a second infinitesimal state-adaptive projective process
$\Psi'_\tau$, with its own limit $\Phi'_t$, that runs the evolution backwards, and prediction and
reconstruction undo each other on pure states, in both orders. In the language of branches, for
every $t > 0$ and every pure state $P$:
\begin{itemize}
\item a predicted branch of $(\Psi_{t/n})^n$ from $P$ has probability tending to one, and a
  reconstructed branch of $(\Psi'_{t/n})^n$ from the predicted state $\Phi_tP$ has probability
  tending to one and ends at $P$;
\item the same with the roles of $\Psi$ and $\Psi'$ exchanged.
\end{itemize}
By P3 this is the same as $\Phi'_t(\Phi_tP) = P$ and $\Phi_t(\Phi'_tP) = P$ for every pure $P$,
which is what the formalization assumes (\lean{DrProj}). The Lean proof then stands as it is. For
the resolvents $\Psi'_\tau = (I + \tau L)^{-1}$, which the brief already names; the brief may also
write it as $\Psi_{-\tau}$, so that DR-PROJ covers both directions of time at once.

\paragraph{Why.}
\begin{itemize}
\item It is what the brief already says. The proof needs ``every pure future state reconstructs a
  unique pure initial state'', and the remark ``The finite process is not unitary'' says that
  ``in the infinitesimal limit \ldots\ prediction and reconstruction become mutually inverse''.
  Mutually inverse means both orders.
\item One order is not enough. $\Phi'_t\circ\Phi_t = I$ on pure states makes $\Phi_t$ injective
  there, but Proposition~17 needs it onto the pure states as well. That comes from
  $\Phi_t\circ\Phi'_t = I$: every pure state is the evolution of the state it reconstructs to.
\item Probabilities alone say nothing about where the reconstruction ends. A branch that is
  certain but ends at a state other than $P$ does not reconstruct $P$, so the reconstructed
  branch must start from the predicted state and end at the initial one.
\item Physically, reconstruction is retrodiction: from the present state one infers the past
  state with certainty. The two sides of the guess are retrodiction of every prediction and
  prediction of every retrodiction.
\end{itemize}

\paragraph{Remark.} In finite dimension reconstruction may not be needed at all. Each $\Psi_\tau$
fixes the maximally mixed state $I/d$, since $\sum_a P_a(I/d)P_a = I/d$, so $\Phi_t$ is affine,
unital and maps pure states to pure states. By the classification of linear maps that preserve
pure states, it is then conjugation by a unitary or an anti-unitary, as the constant maps onto one
pure state are not unital for $d \ge 2$. This sketch is not checked in Lean. In infinite dimension
there is no maximally mixed state, and an isometry that is not unitary is not excluded this way,
so reconstruction is needed there, or at least nothing shown here replaces it.

\paragraph{Confidence.}
\begin{itemize}
\item High: the reconstruction inverts the evolution on pure states, in both orders; the brief
  says ``mutually inverse''.
\item Medium: the wording through branches. The author may prefer to state reconstruction directly
  through the limits, $\Phi'_t\Phi_t = \Phi_t\Phi'_t = I$ on pure states.
\end{itemize}

%======================================================================
\section{Mathematical prerequisites to add to the brief}

Each item is a purely mathematical statement for Lean.

\begin{enumerate}
\renewcommand{\labelenumi}{P\arabic{enumi}.}
\item \textbf{Composition from subdivisions.} Let $(\Psi_\tau)_{\tau > 0}$ be maps of the density
  operators that are affine, and so contract the trace norm. Suppose that for every $t > 0$ the
  ordered products $\Psi_{\tau_k}\circ\dots\circ\Psi_{\tau_1}$ over subdivisions of $[0,t]$
  converge on every density operator in trace norm to $\Phi_t$, as the longest interval tends to
  zero. Then $\Phi_{s+t} = \Phi_t\circ\Phi_s$ for all $s, t > 0$.
  Proof: for subdivisions $A_n$ of $[0,t]$ and $B_n$ of $[0,s]$,
  $\|A_nB_n\rho - \Phi_t\Phi_s\rho\|_1 \le \|B_n\rho - \Phi_s\rho\|_1
  + \|A_n\Phi_s\rho - \Phi_t\Phi_s\rho\|_1 \to 0$, and $A_nB_n$ runs over subdivisions of
  $[0,s+t]$.\newline
  Lean: not formalized. Short, from \lean{AffineOnDensityOperators.traceENorm_sub_le} and
  \lean{isDensityOperator_of_tendsto_traceENorm}.
\item \textbf{Subdivisions of the cutoff resolvents.} Let $U_t = e^{tH/(i\hbar)}$ for a
  self-adjoint, possibly unbounded $H$, with $T_t\rho = U_t\rho U_t^\dagger$, and let $\Psi_\tau$
  be the cutoff resolvents~(23) with $M_\tau$ decreasing in $\tau$. Then
  $\Psi_{\tau_k}\circ\dots\circ\Psi_{\tau_1}\rho \to T_t\rho$ in trace norm over every subdivision
  of $[0,t]$ whose longest interval tends to zero. Proof: if $P_m\rho = \rho$ for the spectral
  projector $P_m$ of $H$ on $[-m,m]$, the states $T_s\rho$ stay in the range of $P_m$, and once
  $M_{\tau_i} \ge m$ each step is the resolvent of the bounded group with generator $H P_m$.
  Telescoping with the bound~(14) gives an error of at most
  $\sum_i 80\,m^2\tau_i^2/\hbar^2 \le 80\,m^2\,t\max_i\tau_i/\hbar^2$. Such $\rho$ are dense, and
  contraction extends the limit to every density operator.\newline
  Lean: not formalized. From \lean{laplaceResolvent_cutoffGroup_eq},
  \lean{traceENorm_laplaceResolvent_sub_unitaryChannel_le},
  \lean{exists_isDensityOperator_mul_eq} and the telescoping bound behind
  \lean{traceENorm_sub_unitaryChannel_le_of_step}.
\item \textbf{A certain branch ends at the limit.} Let $\Psi$ be state-adaptive projective and
  $Q_0, \dots, Q_n$ a branch of $n$ steps of $\Psi$ with probability $w$. Then
  $\|\Psi^nQ_0 - Q_n\|_1 \le 2(1-w)$. Hence, if $(\Psi_{t/n})^n\rho \to \Phi_t\rho$ and branches
  from a pure $P$ have probability tending to one and end at states that converge to $R$, then
  $R = \Phi_tP$. With the reconstruction of the guess this gives $\Phi'_t(\Phi_tP) = P$ and
  $\Phi_t(\Phi'_tP) = P$, so $\Phi_t$ is a bijection of the pure states.\newline
  Lean: the bound is \lean{traceENorm_iterate_sub_le_of_isBranch}; the bijection is
  \lean{DrProj.drEv}. The step from branches to limits is not stated, but follows at once.
\end{enumerate}

\noindent
The proof of Theorem~18 then uses P1 for composition, P2 for the strengthened DR-UNIT
$\Rightarrow$ DR-PROJ, and P3 with Proposition~17 for the unitary implementation.

%======================================================================
\section{Other changes to the brief}

\begin{enumerate}
\item \textbf{DR-PROJ.} Replace ``applying the same process of duration $t/n$, $n$ times, in the
  limit $n \to \infty$'' by ``applying the process over any subdivision of $t$, in the limit where
  the longest interval tends to zero''. Say that $\Phi_0 = I$. Replace ``the probabilities of the
  predicted and reconstructed state sequences tend to one'' by: there is a reconstruction process,
  again infinitesimal and state-adaptive projective, such that from every pure state the predicted
  sequence and the sequence that reconstructs its starting point from its end both have
  probability tending to one, and the same with prediction and reconstruction exchanged.
\item \textbf{Proof of DR-UNIT $\Rightarrow$ DR-PROJ.} State eqs.~(20) and~(24) for subdivisions,
  citing P2; the estimates do not change. For reconstruction, add that the limit of
  $(I + \tau L_\tau)^{-1}$ is $T_{-t}$, which inverts $T_t$, so both orders hold.
\item \textbf{Proof of DR-PROJ $\Rightarrow$ DR-UNIT.} Replace ``Perfect reconstruction gives the
  converse correspondence: every pure future state reconstructs a unique pure initial state'' by a
  citation of P3. Replace the paragraph on common refinements and the remaining interval, from
  ``The same argument gives composition'' to ``$\Phi_t\circ\Phi_s$ for all $t, s \ge 0$'', by a
  citation of P1. Eq.~(29) can stay as an example. Drop ``Therefore $\Phi_0 = I$'', which is now
  part of DR-PROJ.
\item \textbf{Remark ``Mathematical meaning of prediction and reconstruction''.} Align it with the
  new DR-PROJ: the limit is over subdivisions, and the reconstructed branch starts from the
  predicted state and ends at the initial one. The clause that the predicted states converge can
  go, since P3 gives it.
\end{enumerate}

\paragraph{When the blockers close.} Both brief issues close once the brief adopts these
readings, or the alternatives. In Lean, \lean{IsProjectiveLimit} then takes the limit over
subdivisions. With P1 and P2, \lean{IsProjectiveLimit.isEvolution} no longer assumes composition
for irrational ratios, and $\Phi_0 = I$ becomes part of the definition, so DR-UNIT
$\Leftrightarrow$ DR-PROJ holds without the extra clauses of
\lean{drUnit_iff_drProj_and_map_zero_and_map_add}. \lean{DrProj} already matches the guess on
reconstruction; adding the branch form and P3 is short.

\end{document}
